\section{The base functor and the coreflection theorem}
\label{sec:base}

Theorem~\ref{thm:base} attaches a compact Hausdorff space $X_A$ to every object.
This section shows that the assignment is functorial, that it is in fact a
coreflection, and that it exhibits every object as a bundle over its own base.
Everything here is specific to the gauge $\gauge_1$; see
Warning~\ref{warn:vanishing}.

\subsection{Morphisms preserve centres}

\begin{lemma}\label{lem:centre-preserved}
Let $f:A\to B$ be a PB-morphism for the gauge $\gauge_1$.  Then
$f(\Zc(A))\subseteq\Zc(B)$.
\end{lemma}

\begin{proof}
Let $a\in\Zc(A)$.  Then $\gauge_1(a)=0$ by Lemma~\ref{lem:gauge-basic}(e), so
budget-contraction gives $\gauge_1(f(a))\le0$, hence $\ad_{f(a)}=0$, hence
$f(a)\in\Zc(B)$ --- again by Lemma~\ref{lem:gauge-basic}(e), now read in the
other direction.
\end{proof}

\begin{remark}
This is the payoff of the exact identity $\{\gauge_1=0\}=\Zc(A)$.  Note that
the conclusion is genuinely a restriction on arrows and not a triviality: an
ordinary unital contractive homomorphism can certainly move a central element
off the centre, as $\RR^2\hookrightarrow M_2(\RR)$ does.  The morphism law of
Definition~\ref{def:pb-morphism}, introduced to repair the non-functoriality of
the gauge, turns out to buy functoriality of the \emph{base} as well.
\end{remark}

\begin{theorem}[The base is a functor]\label{thm:base-functor}
The assignments $A\mapsto X_A$ and $f\mapsto X_f$ define a functor
\[
  X:\PB^C_{\gauge_1}\longrightarrow\CompHaus\op .
\]
Explicitly, a PB-morphism $f:A\to B$ restricts by
Lemma~\ref{lem:centre-preserved} to a unital homomorphism
$\Zc(f):C(X_A,\RR)\to C(X_B,\RR)$, which by Theorem~\ref{thm:gelfand} is
$X_f^*$ for a unique continuous map $X_f:X_B\to X_A$.  Functoriality is the
functoriality of $\Spec$.
\end{theorem}

\begin{proof}
Lemma~\ref{lem:centre-preserved} gives the restriction, which is unital and
multiplicative; Theorem~\ref{thm:base} identifies source and target with
$C(X_A,\RR)$ and $C(X_B,\RR)$, and Theorem~\ref{thm:gelfand} converts unital
homomorphisms between such algebras into continuous maps of the spectra, in a
way that takes composites to composites and identities to identities.
\end{proof}

\subsection{The commutative part is coreflective}

Write $\CommPB\subseteq\PB^C_{\gauge_1}$ for the full subcategory of
commutative objects.  By Theorem~\ref{thm:gelfand} these are exactly the
algebras $C(X,\RR)$, and the morphisms between them are exactly the unital
homomorphisms --- the budget condition is vacuous, both gauges being zero.
Hence
\[
  \CommPB\;\simeq\;\CompHaus\op .
\]

\begin{theorem}[Coreflection]\label{thm:coreflection}
The centre functor $\Zc:\PB^C_{\gauge_1}\to\CommPB$ is right adjoint to the
inclusion $\iota:\CommPB\hookrightarrow\PB^C_{\gauge_1}$.  That is, for every
commutative object $C$ and every object $A$ there is a bijection, natural in
both variables,
\[
  \Hom_{\PB}\bigl(\iota C,\,A\bigr)\;\cong\;\Hom_{\CommPB}\bigl(C,\,\Zc(A)\bigr).
\]
Equivalently, $\CommPB\simeq\CompHaus\op$ is a coreflective full subcategory of
$\PB^C_{\gauge_1}$, with coreflector $A\mapsto\Zc(A)$.
\end{theorem}

\begin{proof}
First, the inclusion $\Zc(A)\hookrightarrow A$ is a PB-morphism: it is
isometric, and for $z\in\Zc(A)$ both gauges vanish, so the budget condition
holds with equality.  It is therefore a saturated embedding
(Definition~\ref{def:saturated}).

Given a PB-morphism $g:\iota C\to A$ with $C$ commutative, every $c\in C$ has
$\gauge_1^C(c)=0$, so $\gauge_1^A(g(c))=0$ and $g(c)\in\Zc(A)$ by
Lemma~\ref{lem:gauge-basic}(e).  Hence $g$ factors uniquely through
$\Zc(A)\hookrightarrow A$, and the factorisation is a morphism of $\CommPB$
because it is a unital contractive homomorphism between commutative objects.
Conversely, composing a morphism $C\to\Zc(A)$ with the saturated embedding
$\Zc(A)\hookrightarrow A$ gives a PB-morphism $\iota C\to A$.  The two
constructions are mutually inverse, and naturality in $A$ is
Lemma~\ref{lem:centre-preserved} while naturality in $C$ is immediate.
\end{proof}

\begin{corollary}[Points of an object]\label{cor:points}
For every compact Hausdorff $X$ and every object $A$,
\[
  \Hom_{\PB}\bigl(C(X,\RR),\,A\bigr)\;\cong\;C(X_A,X),
\]
the set of continuous maps $X_A\to X$.  In particular $\RR=C(\mathrm{pt},\RR)$
is an initial object, and $\Hom_{\PB}(C(X,\RR),A)$ never sees anything of $A$
beyond its base.
\end{corollary}

\begin{corollary}[Commutative coproducts]\label{cor:comm-coproducts}
The coproduct of $C(X,\RR)$ and $C(Y,\RR)$ in $\PB^C_{\gauge_1}$ exists and
equals $C(X\times Y,\RR)$.
\end{corollary}

\begin{proof}
Let $f:C(X,\RR)\to A$ and $g:C(Y,\RR)\to A$ be PB-morphisms.  By the proof of
Theorem~\ref{thm:coreflection} both land in $\Zc(A)$; in particular their
images commute with one another, so the map $u\otimes v\mapsto f(u)g(v)$ is a
well-defined unital homomorphism on the algebraic tensor product, contractive
for the sup norm, and extends to $C(X\times Y,\RR)=C(X,\RR)\otimes C(Y,\RR)$.
Its image lies in $\Zc(A)$, so its budget vanishes identically and it is a
PB-morphism.  Uniqueness holds because the image of $C(X,\RR)\otimes C(Y,\RR)$
is generated by the two images.
\end{proof}

\begin{warning}[The coproduct is not the free product]
\label{warn:not-free-product}
Corollary~\ref{cor:comm-coproducts} disposes of a suggestion in the earlier
notes, where cocompleteness was to be established by showing that the free
product $A*B$ of PB-algebras is again a PB-algebra.  The free product of
$C(X,\RR)$ and $C(Y,\RR)$ is highly noncommutative, whereas their coproduct in
$\PB$ is the commutative algebra $C(X\times Y,\RR)$.  So the free product is
\emph{not} the coproduct of $\PB$, and the programme ``check that $A*B$ is a
PB-algebra'' does not address cocompleteness.  What Theorem~\ref{thm:coreflection}
shows is that, along the commutative part at least, the coproduct is the one
forced by the geometry: the product of the bases.  The correct general question
is Problem~\ref{prob:coproducts}.
\end{warning}

\subsection{The bundle picture}
\label{sec:bundle}

Theorem~\ref{thm:base} makes $A$ a unital Banach module over $C(X_A,\RR)$, the
action being multiplication by central elements; submultiplicativity gives
$\norm{f\cdot a}\le\norm{f}_\infty\norm{a}$.  This is the standard setting for a
fibrewise decomposition.

\begin{definition}[Fibres]\label{def:fibres}
For $x\in X_A$ let $I_x\subseteq C(X_A,\RR)$ be the ideal of functions
vanishing at $x$, let $J_x=\overline{I_x\cdot A}$ --- a closed two-sided ideal
of $A$, since $I_x\cdot A$ is central-by-$A$ --- and set $A_x=A/J_x$, with
quotient map $\pi_x$.  We call $A_x$ the \emph{fibre of $A$ over $x$}.
\end{definition}

\begin{definition}[Local norm]\label{def:local-norm}
For $a\in A$ and $x\in X_A$ put
\[
  N_x(a)=\inf\bigl\{\norm{f\cdot a}\;:\;f\in C(X_A,\RR),\ 0\le f\le1,\
  f\equiv1\ \text{on a neighbourhood of }x\bigr\}.
\]
\end{definition}

\begin{proposition}\label{prop:usc-field}
For every $a\in A$:
\begin{enumerate}
\item[(a)] $\norm{\pi_x(a)}\le N_x(a)\le\norm{a}$;
\item[(b)] $x\mapsto N_x(a)$ is upper semicontinuous on $X_A$;
\item[(c)] each $\pi_x$ is contractive and budget-contracting, and
  $\pi_x(f\cdot1)=f(x)1$ for $f\in C(X_A,\RR)$.
\end{enumerate}
\end{proposition}

\begin{proof}
(a) If $f\equiv1$ near $x$ then $(1-f)(x)=0$, so $(1-f)\cdot a\in I_x\cdot A\subseteq J_x$
and hence $\pi_x(a)=\pi_x(f\cdot a)$, giving
$\norm{\pi_x(a)}\le\norm{f\cdot a}$; take the infimum.  The upper bound is
$\norm{f\cdot a}\le\norm{f}_\infty\norm{a}\le\norm{a}$.

(b) Given $\varepsilon>0$ choose $f$ admissible for $x$ with
$\norm{f\cdot a}<N_x(a)+\varepsilon$, say $f\equiv1$ on the open set $U\ni x$.
Every $x'\in U$ has $f\equiv1$ on the same neighbourhood $U$ of $x'$, so $f$ is
admissible for $x'$ and $N_{x'}(a)\le\norm{f\cdot a}<N_x(a)+\varepsilon$.

(c) Proposition~\ref{prop:category}, and $f\cdot1-f(x)1\in I_x\cdot A$.
\end{proof}

\begin{remark}[What is missing, precisely]\label{rem:local-norm-gap}
For $C^*$-algebras one has $\norm{\pi_x(a)}=N_x(a)$, and the proof uses an
approximate identity to show that every element of $J_x$ is approximated by
elements $(1-f)\cdot a$.  We do not have that here, and we do not assume it: the
inequality in (a) is all that is claimed.  Consequently
Proposition~\ref{prop:usc-field}(b) gives upper semicontinuity of the
\emph{local} norm $N_x$, and upper semicontinuity of the genuine fibre norm
$x\mapsto\norm{\pi_x(a)}$ follows only once the two agree.
\end{remark}

\begin{question}[Faithfulness of the bundle]\label{q:faithful-bundle}
Is $\norm{a}=\sup_{x\in X_A}\norm{\pi_x(a)}$ for every $a\in A$, and is
$N_x=\norm{\pi_x(\cdot)}$?
\end{question}

For $C^*$-algebras the corresponding statement is a theorem
(Dauns--Hofmann \cite{DaunsHofmann}), and the usual proof needs a
partition-of-unity estimate $\norm{\sum_if_i\cdot a_i}\le\max_i\norm{a_i}$ for
partitions of unity $\{f_i\}$, which holds automatically there but is an extra
hypothesis for a general Banach module --- as is the identification of the
fibre norm with the local norm (Remark~\ref{rem:local-norm-gap}).  A positive answer would give the slogan of this section
its literal meaning:

\begin{quote}
\emph{a PB-algebra is the algebra of sections of an upper semicontinuous field
of pointless PB-algebras over its own classical base.}
\end{quote}

Two ingredients are still missing for the slogan, and both are recorded in
\S\ref{sec:open}: Question~\ref{q:faithful-bundle}, and the question of whether
the fibres are PB-algebras at all, which is the quotient-closure
Problem~\ref{prob:quotients}.  It is worth stressing what is \emph{not}
missing: the base, the central action and the fibring are unconditional, and
they are available because the axioms were built to leave the centre rigid.
